\subsection{The one-dimensional specialization}%
\label{sec:peps_cycle_ground_space}

\begin{definition}[Physical and boundary coordinates of an arc]%
    \label{def:peps_cycle_arc_coordinates}
    \lean{TNLean.PEPS.cycleArcConfigEquiv,
        TNLean.PEPS.arcWord_cycleArcConfigEquiv,
        TNLean.PEPS.cycleArcBoundaryEquiv,
        TNLean.PEPS.cycleArcBoundaryMatrixEquiv,
        TNLean.PEPS.cycleArcPhysicalEquiv}
    \leanok
    \uses{def:cycleTensorOfMPS}
    Let $n\geq3$ and $0<L<n$. Enumerate an arc $R$ of $L$ sites in
    cyclic order. A physical configuration on $R$ is then a word
    $\sigma=(i_1,\ldots,i_L)$. There are exactly two crossing bonds,
    entering and leaving the arc. If a boundary coefficient is
    $x(a,b)$ on their ordered bond values, its boundary matrix is
    $X_{ba}=x(a,b)$. Write $E_R$ for the corresponding ordering of
    physical coefficients.
\end{definition}

\begin{lemma}[Arc contraction and the boundary matrix]%
    \label{lem:peps_cycle_arc_boundary_matrix}
    \lean{TNLean.PEPS.exterior_mul_openRegionMap_cycleTensorOfMPS,
        TNLean.PEPS.cycleArcPhysical_openRegionMap_eq_groundSpaceMap}
    \leanok
    \uses{def:peps_cycle_arc_coordinates,
        def:peps_open_region_contraction, def:ground_space_map}
    Let $A$ be a matrix product tensor of bond dimension $D$, and
    let $\Gamma_R$ be its open contraction on $R$. Let $m$ be the
    number of assignments to the cycle bonds wholly outside $R$.
    Then
    \begin{align}
        m\,E_R\Gamma_R(x)(\sigma)
         & =D^{n-L-1}\operatorname{tr}
        (A^{i_1}\cdots A^{i_L}X).
        \label{eq:peps_cycle_arc_boundary_matrix}
    \end{align}
    If $D>0$, both scalar factors are nonzero, so the regional
    boundary map differs from the MPS boundary map only by an
    invertible scalar change of boundary coordinates.
\end{lemma}

\begin{proof}\leanok
    Summing a regional contraction over all cycle bonds counts it
    once for each exterior assignment. At fixed incoming and
    outgoing bond values $(a,b)$, the same contraction equals
    $D^{n-L-1}(A^{i_1}\cdots A^{i_L})_{ab}$.
    Sum against $x(a,b)$ and use
    $\sum_{a,b}x(a,b)M_{ab}=\operatorname{tr}(MX)$.
    Positive bond dimension permits cancellation of the exterior
    assignment count.
\end{proof}

\begin{theorem}[The regional ground space on a cycle arc]%
    \label{thm:peps_cycle_region_ground_space}
    \lean{TNLean.PEPS.map_regionGroundSpace_cycleTensorOfMPS}
    \leanok
    \uses{def:peps_region_ground_space, def:ground_space_parent,
        def:peps_cycle_arc_coordinates}
    For a matrix product tensor $A$ on a cycle of $n\geq3$ sites,
    and every arc $R$ of length $0<L<n$,
    \begin{align}
        E_R\mathcal G_R
         & =\{\sigma\mapsto
        \operatorname{tr}(A^{i_1}\cdots A^{i_L}X):X\}
        =\mathcal G_L(A).
        \label{eq:peps_cycle_ground_space}
    \end{align}
    No injectivity assumption is needed. This identifies the
    one-dimensional and graph definitions in
    \cite[Section~IV.C.1]{Cirac2021Matrix}, local source lines
    1987--2011.
\end{theorem}

\begin{proof}\leanok
    \uses{lem:peps_cycle_arc_boundary_matrix}
    For $D>0$, the two-bond boundary coefficients range over all
    matrices, and the nonzero scalar changes of boundary coordinates
    in (\ref{eq:peps_cycle_arc_boundary_matrix}) preserve the range.
    For $D=0$, the nonempty arc has an empty boundary-configuration
    space, while the matrix boundary space has only the zero matrix.
    Both ground spaces are therefore zero.
\end{proof}
